% GENERATED FILE. Edit canonical lessons, then run npm run generate:books.
% canonical-source-hash: fnv1a64:a6271bac6b3041f5
% canonical-lessons: MR-C405-viman, MR-C405-uddan, MR-C405-dhavpatti, MR-C405-simasulka, MR-C405-vhisa

\chapter{Planes and Runways}
\label{ch:mr-planes-and-runways}
\hlchaptermodality{405}

\begin{chapteropening}
\textbf{By the end of this chapter:} \emph{I can talk about a plane, a flight, a runway, customs and a visa.}

Complete the last lesson of chapter 405: I can talk about a plane, a flight, a runway, customs and a visa.
\end{chapteropening}

\section[\texorpdfstring{\mr{विमान} (vimān)}{विमान (vimān)}]{\mr{विमान} (vimān) — an aeroplane}

\label{lesson:MR-C405-viman}



\begin{quote}
Before the new one: say the Marathi for to brush against, then the Marathi for to crowd thickly.
\end{quote}

\subsection*{You'll want to know: \mr{विमान}}
\textbf{\mr{विमान}} — \emph{vimān} — \textquotedblleft{}an aeroplane\textquotedblright{}.

\begin{quote}
\mr{आमचं} \mr{विमान} \mr{वेळेवर} \mr{उतरलं}.
\end{quote}

\begin{quote}
\emph{āmcaṁ vimān veḷevar utarlaṁ.}
\end{quote}

\begin{quote}
Our plane landed on time.
\end{quote}

Neuter.

\mr{विमानाने}, by plane; \mr{विमानतळ}, the airport, is built on it.

\begin{cousinweb}
\textbf{\mr{विमान}} (\emph{vimān}) is a tatsama, taken from Sanskrit \emph{vimāna}, which names the flying chariot of the old stories (Monier-Williams).
\end{cousinweb}

\begin{grammarlens}[title={What you've built}]
One more word for getting around.
\end{grammarlens}

\subsection*{Your turn}
\begin{itemize}
\raggedright
  \item \emph{Say it:} \emph{vimān}
  \item \emph{Say it:} \emph{vimān}, once more
  \item \emph{Say it:} \emph{āmcaṁ vimān veḷevar utarlaṁ.} — Our plane landed on time.
  \item \emph{Recall:} say the Marathi for to brush against, then the Marathi for to crowd thickly, then say \emph{vimān} again
\end{itemize}

\begin{tcolorbox}[breakable,colback=teal!4,colframe=teal!35!black,title={Before you move on}]
What does \mr{विमान} mean? (\textquotedblleft{}An aeroplane\textquotedblright{}.) Say it once more.
\end{tcolorbox}
\section[\texorpdfstring{\mr{उड्डाण} (uḍḍāṇ)}{उड्डाण (uḍḍāṇ)}]{\mr{उड्डाण} (uḍḍāṇ) — a flight; a take-off}

\label{lesson:MR-C405-uddan}



\begin{quote}
Before the new one: say the Marathi for to crowd thickly, then the Marathi for an aeroplane.
\end{quote}

\subsection*{You'll want to know: \mr{उड्डाण}}
\textbf{\mr{उड्डाण}} — \emph{uḍḍāṇ} — \textquotedblleft{}a flight; a take-off\textquotedblright{}.

\begin{quote}
\mr{धुक्यामुळे} \mr{सकाळचं} \mr{उड्डाण} \mr{उशिरा} \mr{झालं}.
\end{quote}

\begin{quote}
\emph{dhukyāmuḷe sakāḷcaṁ uḍḍāṇ uśirā jhālaṁ.}
\end{quote}

\begin{quote}
Because of the fog the morning flight was late.
\end{quote}

Neuter.

\mr{उड्डाण} \mr{करणे}, to take off.

\begin{cousinweb}
\textbf{\mr{उड्डाण}} (\emph{uḍḍāṇ}) is built on Sanskrit \emph{uḍḍayana}, \textquotedblleft{}flying up, soaring\textquotedblright{} (Monier-Williams), reshaped in Marathi.
\end{cousinweb}

\begin{grammarlens}[title={What you've built}]
One more word for getting around.
\end{grammarlens}

\subsection*{Your turn}
\begin{itemize}
\raggedright
  \item \emph{Say it:} \emph{uḍḍāṇ}
  \item \emph{Say it:} \emph{uḍḍāṇ}, once more
  \item \emph{Say it:} \emph{dhukyāmuḷe sakāḷcaṁ uḍḍāṇ uśirā jhālaṁ.} — Because of the fog the morning flight was late.
  \item \emph{Recall:} say the Marathi for to crowd thickly, then the Marathi for an aeroplane, then say \emph{uḍḍāṇ} again
\end{itemize}

\begin{tcolorbox}[breakable,colback=teal!4,colframe=teal!35!black,title={Before you move on}]
What does \mr{उड्डाण} mean? (\textquotedblleft{}A flight; a take-off\textquotedblright{}.) Say it once more.
\end{tcolorbox}
\section[\texorpdfstring{\mr{धावपट्टी} (dhāvpaṭṭī)}{धावपट्टी (dhāvpaṭṭī)}]{\mr{धावपट्टी} (dhāvpaṭṭī) — a runway}

\label{lesson:MR-C405-dhavpatti}



\begin{quote}
Before the new one: say the Marathi for an aeroplane, then the Marathi for a flight.
\end{quote}

\subsection*{You'll want to know: \mr{धावपट्टी}}
\textbf{\mr{धावपट्टी}} — \emph{dhāvpaṭṭī} — \textquotedblleft{}a runway\textquotedblright{}.

\begin{quote}
\mr{विमान} \mr{धावपट्टीवर} \mr{उभं} \mr{आहे}.
\end{quote}

\begin{quote}
\emph{vimān dhāvpaṭṭīvar ubhaṁ āhe.}
\end{quote}

\begin{quote}
The plane is standing on the runway.
\end{quote}

Feminine.

\mr{धाव} is a run and \mr{पट्टी} a strip: the strip planes run along.

\begin{cousinweb}
\textbf{\mr{धावपट्टी}} (\emph{dhāvpaṭṭī}) is a Marathi compound of \textbf{\mr{धाव}}, \textquotedblleft{}a run\textquotedblright{}, inherited from \emph{*dhāva} (CDIAL 6800), and \textbf{\mr{पट्टी}}, \textquotedblleft{}a strip\textquotedblright{}.
\end{cousinweb}

\begin{grammarlens}[title={What you've built}]
One more word for getting around.
\end{grammarlens}

\subsection*{Your turn}
\begin{itemize}
\raggedright
  \item \emph{Say it:} \emph{dhāvpaṭṭī}
  \item \emph{Say it:} \emph{dhāvpaṭṭī}, once more
  \item \emph{Say it:} \emph{vimān dhāvpaṭṭīvar ubhaṁ āhe.} — The plane is standing on the runway.
  \item \emph{Recall:} say the Marathi for an aeroplane, then the Marathi for a flight, then say \emph{dhāvpaṭṭī} again
\end{itemize}

\begin{tcolorbox}[breakable,colback=teal!4,colframe=teal!35!black,title={Before you move on}]
What does \mr{धावपट्टी} mean? (\textquotedblleft{}A runway\textquotedblright{}.) Say it once more.
\end{tcolorbox}
\section[\texorpdfstring{\mr{सीमाशुल्क} (sīmāśulka)}{सीमाशुल्क (sīmāśulka)}]{\mr{सीमाशुल्क} (sīmāśulka) — border customs duty}

\label{lesson:MR-C405-simasulka}



\begin{quote}
Before the new one: say the Marathi for a flight, then the Marathi for a runway.
\end{quote}

\subsection*{You'll want to know: \mr{सीमाशुल्क}}
\textbf{\mr{सीमाशुल्क}} — \emph{sīmāśulka} — \textquotedblleft{}border customs duty\textquotedblright{}.

\begin{quote}
\mr{विमानतळावर} \mr{आम्ही} \mr{सीमाशुल्क} \mr{भरलं}.
\end{quote}

\begin{quote}
\emph{vimāntaḷāvar āmhī sīmāśulka bharlaṁ.}
\end{quote}

\begin{quote}
We paid customs duty at the airport.
\end{quote}

Neuter.

The official word on forms and boards; the older word is \mr{जकात}.

\begin{cousinweb}
\textbf{\mr{सीमाशुल्क}} (\emph{sīmāśulka}) is a modern word built from Sanskrit \emph{sīmā}, \textquotedblleft{}boundary\textquotedblright{}, and \emph{śulka}, \textquotedblleft{}toll, price\textquotedblright{} (Monier-Williams): the toll paid at the border.
\end{cousinweb}

\begin{grammarlens}[title={What you've built}]
One more word for getting around.
\end{grammarlens}

\subsection*{Your turn}
\begin{itemize}
\raggedright
  \item \emph{Say it:} \emph{sīmāśulka}
  \item \emph{Say it:} \emph{sīmāśulka}, once more
  \item \emph{Say it:} \emph{vimāntaḷāvar āmhī sīmāśulka bharlaṁ.} — We paid customs duty at the airport.
  \item \emph{Recall:} say the Marathi for a flight, then the Marathi for a runway, then say \emph{sīmāśulka} again
\end{itemize}

\begin{tcolorbox}[breakable,colback=teal!4,colframe=teal!35!black,title={Before you move on}]
What does \mr{सीमाशुल्क} mean? (\textquotedblleft{}Border customs duty\textquotedblright{}.) Say it once more.
\end{tcolorbox}
\section[\texorpdfstring{\mr{व्हिसा} (vhisā)}{व्हिसा (vhisā)}]{\mr{व्हिसा} (vhisā) — a visa}

\label{lesson:MR-C405-vhisa}



\begin{quote}
Before the new one: say the Marathi for a runway, then the Marathi for border customs duty.
\end{quote}

\subsection*{You'll want to know: \mr{व्हिसा}}
\textbf{\mr{व्हिसा}} — \emph{vhisā} — \textquotedblleft{}a visa\textquotedblright{}.

\begin{quote}
\mr{मीराला} \mr{व्हिसा} \mr{मिळाला}.
\end{quote}

\begin{quote}
\emph{mīrālā vhisā miḷālā.}
\end{quote}

\begin{quote}
Meera got her visa.
\end{quote}

Masculine.

\mr{व्हिसा} \mr{मिळणे} takes the dative: \mr{तिला} \mr{व्हिसा} \mr{मिळाला}.

\begin{cousinweb}
\textbf{\mr{व्हिसा}} (\emph{vhisā}) is a loan from English \emph{visa}.
\end{cousinweb}

\begin{grammarlens}[title={What you've built}]
One more word for getting around.
\end{grammarlens}

\subsection*{Your turn}
\begin{itemize}
\raggedright
  \item \emph{Say it:} \emph{vhisā}
  \item \emph{Say it:} \emph{vhisā}, once more
  \item \emph{Say it:} \emph{mīrālā vhisā miḷālā.} — Meera got her visa.
  \item \emph{Recall:} say the Marathi for a runway, then the Marathi for border customs duty, then say \emph{vhisā} again
\end{itemize}

\begin{tcolorbox}[breakable,colback=teal!4,colframe=teal!35!black,title={Before you move on}]
What does \mr{व्हिसा} mean? (\textquotedblleft{}A visa\textquotedblright{}.) Say it once more.
\end{tcolorbox}
